\section{Funciones meromorfas}

\subsection{Funciones meromorfas}

\begin{definición} [Función meromorfa]
    Sea $\Omega \subset \mathbb{C}$ un abierto y una función $f: \Omega \to \mathbb{C}$. Entonces diremos que $f$ es \textbf{meromorfa} en $\Omega$ si es holomorfa en $\Omega \setminus S$, para algún conjunto $S \subset \Omega$ sin puntos de acumulación, y tal que cada $z \in S$ es un polo de $f$.
\end{definición}

\begin{observación}
    Una función $f$ es meromorfa en $\Omega$ si $f$ es holomorfa en $\Omega$ salvo en un conjunto de puntos aislados de $\Omega$ en los que $f$ tiene polos.
\end{observación}

\ejemplo{
    \begin{itemize}
        \item Toda función holomorfa en $\Omega$ es meromorfa en $\Omega$.
        \item Si $p$ y $q$ son polinomios, $q \not \equiv 0$, entonces la función racional $f(z) = \frac{p(z)}{q(z)}$ es meromorfa en $\mathbb{C}$.
        \item Las funciones $\sin(z)$, $\cos(z)$, $\tan(z)$, $\cot(z)$, $\sec(z)$ y $\csc(z)$ son meromorfas en $\mathbb{C}$.
        \item Si $f(z) = \frac{g(z)}{h(z)}$ con $g$ y $h$ funciones holomorfas en $\Omega$ y $h \not \equiv 0$, entonces $f$ es meromorfa en $\Omega$.
        \item Las funciones $\sin(\frac{1}{z})$ y $e^{\frac{1}{z}}$ son meromorfas en $\mathbb{C} \setminus \{0\}$, pero no en $\mathbb{C}$, pues en $z=0$ tienen una singularidad esencial.
    \end{itemize}
}

Aunque no veremos la demostración en este curso, toda función meromorfa en un dominio $\Omega$ puede escribirse localmente como el cociente de dos funciones holomorfas. Este resultado es una consecuencia del teorema de factorización de Weierstrass para funciones holomorfas.

Si $f : \Omega \to \mathbb{C}$ es meromorfa en $\Omega$, sabemos que $f \in \mathcal{H}(\Omega \setminus S)$, donde $S$ es el conjunto de polos de $f$. Así, podemos considerar
\[
    \begin{matrix}
        f : \Omega &\to & \mathbb{C}_\infty \\
        z &\mapsto & \begin{cases}
            f(z), & z \in \Omega \setminus S \\[6pt]
            \infty, & z \in S
        \end{cases}
    \end{matrix}
\]
De modo que $f$ es holomorfa en $\Omega \setminus S$ y continua en $\Omega$ ($f: \Omega \to \mathbb{C}_\infty$).

Denotamos por
\[
    \mathcal{M}(\Omega) = \{ f : \Omega \to \mathbb{C}_\infty : f \text{ es meromorfa en } \Omega \}
\]
Tenemos entonces que $\mathcal{H}(\Omega) \subset \mathcal{M}(\Omega)$.
\[
    f \in \mathcal{M}(\Omega) \implies \frac{1}{f} \in \mathcal{M}(\Omega)
\]  
y se tiene que $\mathcal{M}(\Omega)$ es un cuerpo con las operaciones de suma y multiplicación de funciones, mientras que $\mathcal{H}(\Omega)$ tiene estructura algebraica de anillo (pero NO de cuerpo). Además, $f \in \mathcal{M}(\Omega) \implies f' \in \mathcal{M}(\Omega)$.

\subsection{Principio del argumento}

\begin{teorema} [Principio del argumento]
    Sea $\Omega \subset \mathbb{C}$ un dominio y sea $f \in \mathcal{M}(\Omega)$. Sea $\{z_j : j \in J\}$ con $J \subset \mathbb{N}$ el conjunto de ceros de $f$ en $\Omega$ repetidos tantas veces como indica su multiplicidad, y sea $\{p_k : k \in \tilde{J}\}$ con $\tilde{J} \subset \mathbb{N}$ el conjunto de polos de $f$ en $\Omega$ repetidos tantas veces como indica su orden. Sea $\gamma$ una curva cerrada homótopa a $0$ en $\Omega$ que no pasa por ningún cero ni por ningún polo de $f$. Entonces,
    \[
        \frac{1}{2 \pi i} \int_{\gamma} \frac{f'(z)}{f(z)} \, dz = \sum_{j \in J} \Ind_{\gamma}(z_j) - \sum_{k \in \tilde{J}} \Ind_{\gamma}(p_k) = \Ind_{f \circ \gamma}(0)
    \] 
\end{teorema}

\begin{proof}
    \[
        \frac{1}{2 \pi i} \int_{\gamma} \frac{f'(z)}{f(z)} \, dz = \frac{1}{2 \pi i} \int_{a}^{b} \frac{f'(\gamma(t))}{f(\gamma(t))} \gamma'(t) \, dt = \frac{1}{2 \pi i} \int_{a}^{b} \frac{(f \circ \gamma)'(t)}{(f \circ \gamma)(t)} \, dt = \frac{1}{2 \pi i} \int_{f \circ \gamma} \frac{1}{w - 0} \, dw = \Ind_{f \circ \gamma}(0)
    \]
    Para probar la primera igualdad usando el lema del Tema 12 (usado también en el Tema 14), restringiéndonos a $\Omega_1$, tenemos que hay una cantidad finita de ceros y polos de $f$ en $\Omega_1$, y que $\Ind_{\gamma}(z) = 0$ para todo $z \in \Omega \setminus \overline{\Omega}_1$. Con lo cual,
    \[
        \sum_{j \in J} \Ind_{\gamma}(z_j) - \sum_{k \in \tilde{J}} \Ind_{\gamma}(p_k) = \sum_{\substack{j \in J \\ z_j \in \Omega_1}} \Ind_{\gamma}(z_j) - \sum_{\substack{k \in \tilde{J} \\ p_k \in \Omega_1}} \Ind_{\gamma}(p_k)
    \]
    Sean $z_1, \ldots, z_m$ los ceros de $f$ en $\Omega_1$ y $p_1, \ldots, p_l$ los polos de $f$ en $\Omega_1$, entonces como $f \in \mathcal{M}(\Omega)$ (y por tanto $f \in \mathcal{M}(\Omega_1)$) tenemos que $f \in \mathcal{H}(\Omega_1 \setminus \{p_1, \ldots, p_l\})$. Por la \cref{prop:ceros_aislados} tenemos que $\exists g \in \mathcal{H}(\Omega_1 \setminus \{p_1, \ldots, p_l\})$ tal que
    \[
        f(z) = (z - z_1) \cdots (z - z_m) g(z)  \quad \forall z \in \Omega_1 \setminus \{p_1, \ldots, p_l\}
    \]  
    con $g(z) \neq 0$ para todo $z \in \Omega_1$. 
    Observamos que los polos de $g$ son los polos de $f$. Entonces tenemos por la \cref{prop:caracterizacion_polos} que $\exists h \in \mathcal{H}(\Omega_1)$ tal que
    \[
        g(z) = \frac{h(z)}{(z - p_1)(z - p_2) \cdots (z - p_l)} \quad \forall z \in \Omega_1 \setminus \{p_1, \ldots, p_l\}
    \]
    y $h$ no se anula en $\Omega_1$ (porque $g$ no se anula en $\Omega_1$). Por tanto,
    \[
        f(z) = (z-z_1) \cdots (z-z_m) \frac{h(z)}{(z - p_1)(z - p_2) \cdots (z - p_l)} \quad \forall z \in \Omega_1 \setminus \{p_1, \ldots, p_l\}
    \]
    Puesto que $\frac{f'}{f} = (\log f)'$, derivando obtenemos
    \[
        \frac{f'(z)}{f(z)} = \left( \log \left( (z-z_1) \cdots (z-z_m) \frac{h(z)}{(z - p_1)(z - p_2) \cdots (z - p_l)} \right) \right)'
    \]
    \[
        = \left( \log \left( \sum_{j=1}^{m} (z - z_j) \right) + \log(h(z)) - \log \left( \sum_{k=1}^{l} (z - p_k) \right) \right)' 
    \]
    \[
        = \left( \sum_{j=1}^{m} \log(z - z_j) \right)' + \left( \log(h(z)) \right)' - \left( \sum_{k=1}^{l} \log(z - p_k) \right)' 
        = \sum_{j=1}^{m} \frac{1}{z - z_j} + \frac{h'(z)}{h(z)} - \sum_{k=1}^{l} \frac{1}{z - p_k}
    \]
    
    Como $h \in \mathcal{H}(\Omega_1)$ y no se anula en $\Omega_1$ entonces tenemos que $\frac{h'(z)}{h(z)}$ es holomorfa en $\Omega_1$ y puesto que $\gamma$ es homótopa a $0$ en $\Omega_1$, por el teorema de Cauchy,
    \[
        \int_{\gamma} \frac{h'(z)}{h(z)} \, dz = 0
    \]
    Así, integrando
    \[
        \frac{1}{2 \pi i} \int_{\gamma} \frac{f'(z)}{f(z)} \, dz = \sum_{j=1}^{m} \frac{1}{2 \pi i} \int_{\gamma} \frac{1}{z - z_j} \, dz - \sum_{k=1}^{l} \frac{1}{2 \pi i} \int_{\gamma} \frac{1}{z - p_k} \, dz = \sum_{j=1}^{m} \Ind_{\gamma}(z_j) - \sum_{k=1}^{l} \Ind_{\gamma}(p_k)
    \]
\end{proof}

\begin{corolario}
    Algo interesante de notar, es que si $\gamma$ es una curva simple, se cumple que $\Ind_\gamma(z) \in \{0,1\} \forall z \in \mathbb{C} \setminus \gamma^*$.\\
    Aplicando ahora el principio del argumento, obtenemos que si $f \in \mathcal{M}(\Omega)$ y $\gamma$ es una curva simple homótopa a $0$ en $\Omega$ que no pasa por ningún cero ni por ningún polo de $f$, entonces
    \[
        \frac{1}{2 \pi i} \int_{\gamma} \frac{f'(z)}{f(z)} \, dz = \text{nº de ceros de } f \text{ en el interior de } \gamma - \text{nº de polos de } f \text{ en el interior de } \gamma
    \]
    Esto ultimo tambien se cumple unicamente con que $\Ind_{\gamma}(z) \in \{0,1\} \forall z \in \mathbb{C} \setminus \gamma^*$, es decir, no es necesario que $\gamma$ sea simple.
\end{corolario}

El teorema anterior se llama principio del argumento ya que en $\gamma$ podemos definir una cantidad finita de discos, todos ellos con centro en un punto de la imagen de $f \circ \gamma$ y radio constante entre ellos, de modo que al recorrer la curva $\gamma$ en sentido positivo, sin pasar por ninguno de los ceros ni polos de $f$, es decir:
\[
\exists \{t_0, t_1, \ldots, t_m\}, \delta >0 :
\gamma ([t_{k-1}, t_k]) \subset D(f(\gamma(t_k)), \delta) \quad \forall k = 1, \ldots, m
\]
Dejandonos así la siguiente igualdad:
\[
\int_{\gamma_j} \frac{f'(z)}{f(z)} \, dz = 
log_j(f(\gamma(t_j))) - log_j(f(\gamma(t_{j-1})))
\]
Siendo cada $log_j$ una determinación holomorfa del logaritmo en cada disco $D(f(\gamma(t_j)), \delta)$, y siendo $\gamma_j = \gamma|_{[t_{j-1}, t_j]}$.\\
Así, sumando todas las integrales, obtenemos que:
\[
\int_{\gamma} \frac{f'}{f} \, dz = 
\sum_{j=1}^{m} \int_{\gamma_j} \frac{f'(z)}{f(z)} \, dz = \
\sum_{j=1}^{m} \left( log_j(f(\gamma(t_j))) - log_j(f(\gamma(t_{j-1}))) \right) =
log_m(f(\gamma(t_m))) - log_1(f(\gamma(t_0))) =
\]
\[
= \left( \text{arg}_m(f(\gamma(t_m))) - \text{arg}_0(f(\gamma(t_0))) \right) + i \left( \log|f(\gamma(t_m))| - \log|f(\gamma(t_0))| \right) =
2 \pi i N
\]

Siendo $\text{arg}_m(f(\gamma(t_m))) - \text{arg}_0(f(\gamma(t_0))) = 2 \pi N$.

Veamos ahora un resultado muy útil que se deduce del principio del argumento.

\subsection{Teorema de Rouché y Teorema de Hurwitz}

\begin{teorema}[Teorema de Rouché]
    Sea $\Omega \subset \mathbb{C}$ un dominio y sean $f, g$ meromorfas en $\Omega$. Sea $\gamma$ una curva cerrada homótopa a $0$ en $\Omega$ que cumple que $\Ind_{\gamma}(z) \in \{0,1\} \ \forall z \in \Omega \setminus \gamma^*$.\\
    Si $|f(z) - g(z)| < |f(z)|$ para todo $z \in \gamma^*$, entonces $f$ y $g$ cumplen:
    \[
        \text{nº de ceros de } f \text{ en el interior de } \gamma - \text{nº de polos de } f \text{ en el interior de } \gamma =
    \]
    \[
    \text{nº de ceros de } g \text{ en el interior de } \gamma - \text{nº de polos de } g \text{ en el interior de } \gamma
    \]
\end{teorema}

\begin{proof}
    Por hipótesis, $f$ no se anula en $\gamma^*$, y $g$ tampoco se anula en $\gamma^*$. Además, $g$ no tiene polos en $\gamma^*$.\\
    Tomemos ahora $h(z) = \frac{g(z)}{f(z)} \quad \forall z \in \Omega$. Entonces $h$ es meromorfa en $\Omega$ y no tiene ceros ni polos en $\gamma^*$.\\
    Ahora tomemos $z \in \gamma^*$. Entonces:
    \[
    \left| h(z) - 1 \right| = 
    \left| \frac{g(z)}{f(z)} - 1 \right| =
    \frac{|g(z) - f(z)|}{|f(z)|} <
    \frac{|f(z)|}{|f(z)|} = 1
    \]
    Esto nos dice que $h(\gamma^*) \subset D(1,1)$, luego $0$ pertenece a la componente conexa no acotada de $h \circ \gamma$, ya que $0 \in \mathbb{C} \setminus D(1,1 - \delta) \subset \mathbb{C} \setminus h(\gamma^*)$ para algún $\delta > 0$.\\
    Por tanto, $\Ind_{h \circ \gamma}(0) = 0$. Usando el principio del argumento para $h$, tenemos:
    \[
    0 =
    \Ind_{h \circ \gamma}(0) \underbrace{=}_{\text{Ppio. argumento}}
    \frac{1}{2 \pi i} \int_{\gamma} \frac{h'(z)}{h(z)} \, dz =
    \frac{1}{2 \pi i} \int_{\gamma} \frac{1}{h(z)} \left( \frac{f(z)g'(z) - f'(z)g(z)}{\left(f(z)\right)^2} \right) \, dz =
    \]
    \[
    = \frac{1}{2 \pi i} \int_{\gamma} \frac{f(z)}{g(z)} \cdot \left( \frac{f(z)g'(z) - f'(z)g(z)}{\left(f(z)\right)^2} \right) \, dz =
    \frac{1}{2 \pi i} \left( \int_{\gamma} \frac{g'(z)}{g(z)} \, dz - \int_{\gamma} \frac{f'(z)}{f(z)} \, dz \right) =
    \]
    \[
    \underbrace{=}_{\Ind_{\gamma}(z) \in \{0,1\}}
    \left( \text{nº de ceros de } g \text{ en el interior de } \gamma - \text{nº de polos de } g \text{ en el interior de } \gamma \right) -
    \]
    \[
    \left( \text{nº de ceros de } f \text{ en el interior de } \gamma - \text{nº de polos de } f \text{ en el interior de } \gamma \right)
    \]
    Teniendo así la igualdad buscada.
\end{proof}

\ejemplo{
    Veamos cuantos ceros tiene $p(z) = 2z^10 + 4z^2 +1$ en el disco unitario $D(0,1)$.\\
    Desde luego, $p \in \mathcal{M}(\mathbb{C})$. Tomemos entonces $f(z) = 4z^2$. Si $z \in C(0,1)$, entonces:
    \[
        |f(z)| = 4|z|^2 = 4 |z|^2 = 4 \cdot 1^2 = 4
    \]
    Luego:
    \[
        |f(z) - p(z)| = 
        \left| \left( 4z^2 \right) - \left( 2z^{10} + 4z^2 + 1 \right) \right| = 
        \left| -2z^{10} - 1 \right|
    \] 
    \[
    \leq
        2|z|^{10} + 1 =
        2 \cdot 1^{10} + 1 = 3 < 4 = |f(z)| \quad \forall z \in C(0,1)
    \]
    Por el teorema de Rouché, $p$ y $f$ tienen el mismo número de ceros en el interior de $C(0,1)$.\\
    Pero $f(z) = 4z^2$ tiene 2 ceros en $D(0,1)$, por lo que $p$ también tiene 2 ceros en $D(0,1)$ según el teorema de Rouché.\\
    Así, como $p(z)$ tiene grado 10, $p$ tiene 8 ceros fuera del disco unitario, estando todos ellos en el anillo $A(0,1,2)$.
}

\ejemplo{
    Si $\Omega \subset \mathbb{C}$ es un dominio y $\Omega \supset \overline{D} = \overline{D(0,1)}$, y si $f \in \mathcal{H}(\Omega)$ cumple que $f\left(C(0,1)\right) \subset D(0,1)$, entonces $f$ tiene un único punto fijo en $D(0,1)$.\\
    Notemos que $z_0$ es un punto fijo de $f$ si y solo si $f(z_0) - z_0 = 0$, o lo que es lo mismo, si es un cero de $g(z) = f(z) - z$.\\
    Sea ahora $h(z) = z \quad \forall z \in \mathbb{C}$, entonces:
    \[
    \left| g(z) + h(z) \right| =
    \left| g(z) - (-h(z)) \right| =
    \left| f(z) - z + z \right| =
    |f(z)| < 
    1 = 
    \left| -h(z) \right| =
    \left| -z \right| \quad \forall z \in C(0,1)
    \]
    Luego $g$ y $-h$ tienen el mismo número de ceros en el interior de $C(0,1)$ por el teorema de Rouché, al no tener polos, ya que $h$ tiene 1 unico cero en $D(0,1)$ (de multiplicidad 1).\\

}

Veamos a continuación otra aplicación del teorema de Rouché en el siguiente resultado.

\begin{teorema} [Teorema de Hurwitz]
    Sea $\Omega \subset \mathbb{C}$ un dominio y sea $(f_n)_{n \in \mathbb{N}}$ una sucesión de funciones holomorfas en $\Omega$, que converge uniformemente sobre compactos de $\Omega$ a una función $f$. Si para cada $n \in \mathbb{N}$, $f_n$ no se anula en $\Omega$, entonces o bien $f$ es identicamente nula en $\Omega$, o bien $f$ no se anula en $\Omega$.
\end{teorema}

\textbf{Nota:} Puede ocurrir que $f \equiv 0$. Por ejemplo, si tomamos $f_n = \frac{1}{n} \quad \forall n \in \mathbb{N}$, entonces $f_n \to f \equiv 0$ uniformemente sobre compactos de $\mathbb{C}$.

\begin{proof}
    Supongamos que $f \not \equiv 0$ y supongamos que existe $z_0 \in \Omega$ tal que $f(z_0) = 0$. Entonces, como $f \not \equiv 0$, tenemos que $z_0$ es un cero aislado de $f$. Sea entonces $r > 0$ tal que $\overline{D}(z_0, r) \subset \Omega$ y $f(z) \neq 0 \quad \forall z \in D'(z_0, r)$.\\
    Por lo tanto, como $f$ no se anula en $C := C(z_0, \frac{r}{2})$, tenemos que $\min \{|f(z)| : z \in C\} = \delta > 0$.\\
    Como $f_n \to f$ uniformemente sobre compactos de $\Omega$, entonces como $C$ es compacto, existe $N \in \mathbb{N}$ tal que:
    \[
        \sup_{z \in C} |f_n(z) - f(z)| < \delta \quad \forall n \geq N
    \]
    Como $|f_N(z) - f(z)| < \delta \leq |f(z)|$ para todo $z \in C$, por el teorema de Rouché, deducimos que $f_N$ tiene en $D(z_0, \frac{r}{2})$ el mismo número de ceros que $f$ en dicho disco, o sea, $1$. Esta contradicción termina la demostración. 
\end{proof}



